\chapter{The Hyperbolic Plane and the Modular Group}

\begin{center}\itshape
A continued fraction is a sequence of integers. A geodesic is a straight path on a curved surface. In the hyperbolic plane, one can draw a tiling whose triangles turn the path into the sequence: count how many triangles it crosses before turning the other way.
\end{center}

Henri Poincaré found an unexpected use for non-Euclidean geometry while studying functions. In 1880, teaching at Caen, he was investigating functions on the upper half-plane unchanged by certain groups of transformations $z\mapsto(az+b)/(cz+d)$. He called them Fuchsian functions, after Lazarus Fuchs. In his later account in \emph{Science et méthode}, he recalled spending two weeks trying to prove that such functions could not exist, then finding a family of them. He left Caen on a geological excursion organized by the School of Mines.

At Coutances, as he stepped onto an omnibus, he recognized the transformations from his work as transformations of non-Euclidean geometry. He continued a conversation without checking the idea, ``with perfect certainty'', and verified it after returning to Caen. The recollection is about a connection between two subjects: the matrices acting on functions also move a geometric plane without changing distances.

We can work out that connection explicitly. The orientation-preserving symmetries of the hyperbolic plane form $PSL_2(\R)$. A matrix records both a point and a direction there; multiplying it by a diagonal matrix follows a geodesic. The integer matrices of Volume~6 produce a tiling, and Caroline Series' cutting-sequence theorem explains how a path through that tiling records a continued fraction. Appendix~\ref{geometry:ch:noneuclid} gives the earlier history of non-Euclidean geometry and its comparison with spherical geometry.

\section{The half-plane model}

In the Euclidean plane, a segment of length $1$ has the same length wherever we put it. Here its height matters. On
\[
\mathbb H=\{z=x+iy\in\mathbb C:y>0\},
\]
a small displacement $dz$ has length $|dz|/y$. At height $2$, the same Euclidean displacement costs half as much as it does at height $1$. Approaching height $0$ along a vertical line requires the divergent integral $\int_0^1dy/y$.

This is a metric on a plane, rather than a surface drawn in three-dimensional space. Hilbert's theorem of 1901 rules out a complete smooth surface of constant curvature $-1$ immersed in $\R^3$; the intrinsic model needs no such realization. Example~\ref{ex:halfplane} computed its curvature directly.

\begin{definition}
For a piecewise $C^1$ path $\gamma:[0,1]\to\mathbb H$, its \emph{hyperbolic length} is
\[
L(\gamma)=\int_0^1\frac{|\gamma'(t)|}{\operatorname{Im}\gamma(t)}\,dt.
\]
The \emph{hyperbolic distance} $d(z,w)$ is the infimum of these lengths over paths joining $z$ to $w$. A measurable region $R$ has \emph{hyperbolic area}
$\iint_R dx\,dy/y^2$.
\end{definition}

At a fixed point, the metric multiplies every vector's Euclidean length by the same positive factor. It therefore leaves angles unchanged. We can measure an angle in a picture using an ordinary protractor, even though a ruler would give the wrong distances.

\begin{proposition}[isometries]\label{geometry:prop:isometries}
If $a,b,c,d$ are real and $ad-bc=1$, the map $g(z)=(az+b)/(cz+d)$ is a bijection of $\mathbb H$ preserving hyperbolic lengths, distances and areas. The reflection $z\mapsto-\bar z$ has the same preservation properties.
\end{proposition}

\begin{proof}
The denominator cannot vanish in $\mathbb H$: when $c\ne0$ its zero is real, and when $c=0$ the determinant condition gives $d\ne0$. Multiplying numerator and denominator by $c\bar z+d$ gives
\[
\operatorname{Im}g(z)=\frac{\operatorname{Im}z}{|cz+d|^2},
\qquad g'(z)=\frac1{(cz+d)^2}.
\]
The first expression is positive. The inverse comes from the real determinant-one matrix $\bigl(\begin{smallmatrix}d&-b\\-c&a\end{smallmatrix}\bigr)$ and also maps $\mathbb H$ into itself, so $g$ is onto. For every path, the factor $|c\gamma+d|^{-2}$ occurs in both $|(g\circ\gamma)'|$ and $\operatorname{Im}(g\circ\gamma)$, and cancels in the length integral. Applying the inverse shows that the two sets of path lengths, hence their infima, agree. For area, the absolute Jacobian is $|g'|^2=|cz+d|^{-4}$; this cancels the change in the factor $y^{-2}$. Reflection keeps $y$ and $|dz|$ fixed and has absolute Jacobian $1$.
\end{proof}

The matrices form $SL_2(\R)$. The matrices $g$ and $-g$ induce the same map, so their transformations form $PSL_2(\R)=SL_2(\R)/\{\pm I\}$. We shall use the same letter for a matrix and its transformation when the meaning is clear.

The candidates for straight lines are the vertical half-lines and the upper half-circles with centres on the real axis. Call them \emph{hyperbolic lines}, or \emph{h-lines}. The transformations just found preserve this family. Here is a local verification, so we need no circle geometry from a later appendix. Translations and positive dilations preserve it. For $J(z)=-1/z$, an equation
\[
\alpha(x^2+y^2)+\beta x+\gamma=0
\]
becomes $\gamma(u^2+v^2)-\beta u+\alpha=0$ after substituting $z=-1/(u+iv)$ and clearing $u^2+v^2$. This is again a vertical line or a circle centred on the real axis. Every determinant-one map with $c\ne0$ is a composition of these transformations, since
\[
\frac{az+b}{cz+d}=\frac ac-\frac1{c^2(z+d/c)}.
\]
For $c=0$ it is a translation followed by a positive dilation. Inverses have the same property, so an h-line maps onto an h-line.

\begin{lemma}\label{geometry:lem:hlines}
Two distinct points $z,w\in\mathbb H$ lie on exactly one h-line. A transformation from Proposition~\ref{geometry:prop:isometries} sends that line to the imaginary axis and sends $z,w$ to $i,it$, respectively, for some $t>1$.
\end{lemma}

\begin{proof}
Write $z=x+iy$ and $w=u+iv$. If $x=u$, their distinct heights prevent any real-centred circle from containing both, so the vertical line is the only choice. If $x\ne u$, a possible centre $c\in\R$ must satisfy $|z-c|^2=|w-c|^2$. Solving gives the unique value
\[
c=\frac{|w|^2-|z|^2}{2(u-x)}.
\]
The circle with this centre through $z$ also contains $w$.

For finite endpoints $a<b$ of this circle, use $h(\zeta)=(\zeta-a)/(b-\zeta)$. Its matrix has positive determinant $b-a$; dividing the matrix by $\sqrt{b-a}$ puts it in $SL_2(\R)$. It sends the endpoints to $0,\infty$, and hence the h-line to the imaginary axis. For a vertical line $x=a$, use $h(\zeta)=\zeta-a$. Now $h(z)=is$ and $h(w)=is'$ for positive, distinct $s,s'$. Divide the image by $s$. If $s'/s<1$, apply $J$, which fixes $i$ and replaces $i(s'/s)$ by $i(s/s')$.
\end{proof}

\begin{theorem}[hyperbolic distance]\label{geometry:thm:hypdist}
For $z,w\in\mathbb H$,
\[
\cosh d(z,w)=1+\frac{|z-w|^2}{2\operatorname{Im}z\operatorname{Im}w}.
\]
For distinct endpoints, the h-line arc joining them is the unique shortest geometric path. Its monotone parametrizations may have different speeds or pauses. For equal endpoints, the shortest path is constant.
\end{theorem}

\begin{proof}
Start with $i,it$, where $t>1$. Write a joining path as $x(s)+iy(s)$. Then
\[
L(\gamma)\ge\int_0^1\frac{|y'(s)|}{y(s)}\,ds
\ge\int_0^1\frac{y'(s)}{y(s)}\,ds=\log t.
\]
The first inequality is an equality precisely when $x'=0$ on each smooth piece; the second requires $y'\ge0$. Thus only the upward vertical segment can minimize length, and its length is $\log t$. The proposed formula gives
$1+(t-1)^2/(2t)=(t+t^{-1})/2=\cosh(\log t)$.

For distinct general endpoints, apply the preceding lemma. Distance is unchanged. The expression in the formula is unchanged as well, because
\[
g(z)-g(w)=\frac{z-w}{(cz+d)(cw+d)},
\qquad \operatorname{Im}g(z)=\frac{\operatorname{Im}z}{|cz+d|^2}.
\]
The denominator factors cancel. H-lines and minimizing paths map to h-lines and minimizing paths, so uniqueness follows from the vertical case. If $z=w$, both sides of the formula are $1$; a path has length zero only when it is constant.
\end{proof}

\begin{figure}[ht]
\centering
\fig{geometry-ch10-halfplane}
\caption{Left: the h-line $x^2+y^2=1$ (charcoal) and three h-lines through $P$ that never meet it (rust). Right: the region above the unit circle between
$x=\cos(\pi-\alpha)$ and $x=\cos\beta$ is a triangle with angles $\alpha$, $\beta$ and $0$ (its third vertex is at infinity).}
\label{geometry:fig:halfplane}
\end{figure}

\begin{proposition}[many parallels]\label{geometry:prop:parallels}
For an h-line $\ell$ and a point $P\notin\ell$, infinitely many h-lines through $P$ are disjoint from $\ell$.
\end{proposition}

\begin{proof}
Move $\ell$ to the imaginary axis and, if needed, reflect so that $P=x+iy$ has $x>0$. The circle centred at $c>0$ through $P$ has radius $r=\sqrt{(x-c)^2+y^2}$. Its whole upper half lies to the right of the axis when $c-r>0$, or equivalently
\[
c>\frac{x^2+y^2}{2x}.
\]
Every such centre gives a different h-line through $P$. The vertical line $\operatorname{Re}z=x$ is another one.
\end{proof}

\section{The area of a triangle}

Choose three distinct vertices, not all on one h-line, and join them by h-line arcs. The resulting convex region is a \emph{hyperbolic triangle}. A vertex may also lie on $\R\cup\{\infty\}$; it is then an \emph{ideal vertex}. Its sides approach the boundary without reaching it at finite distance, and its angle is taken to be $0$. We exclude collapsed triangles, whose angles would need a separate convention.

\begin{think}
Can all three angles be $60^\circ$? Can all three be $0^\circ$? Try to decide before computing an area.
\end{think}

\begin{theorem}[Gauss--Bonnet for triangles]\label{geometry:thm:hyparea}
A hyperbolic triangle with angles $\alpha,\beta,\gamma$ has area $\pi-\alpha-\beta-\gamma$.
\end{theorem}

\begin{proof}
First suppose that one vertex is ideal. Move it to $\infty$; the two sides ending there become vertical lines. Translate and dilate so that the circle carrying the third side is the unit circle. The region is
\[
x_0\le x\le x_1,\qquad y\ge\sqrt{1-x^2}.
\]
Angles between tangents equal angles between their perpendiculars. Thus, with left and right angles $\alpha,\beta$, its lower endpoints have $x_0=\cos(\pi-\alpha)$ and $x_1=\cos\beta$. Integrating first in $y$ gives
\[
\operatorname{area}=\int_{x_0}^{x_1}\frac{dx}{\sqrt{1-x^2}}
=\arccos x_0-\arccos x_1=\pi-\alpha-\beta.
\]
An ideal endpoint of the lower arc gives $x_0=-1$ or $x_1=1$. The same convergent improper integral covers triangles with two or three ideal vertices.

For a triangle $ABC$ with finite vertices, extend the side from $A$ through $C$ to an ideal point $D$. The larger triangle $ABD$ consists of $ABC$ and $CBD$, sharing only $BC$. Let $\delta$ be the angle of $CBD$ at $B$. The angle of $ABD$ there is $\beta+\delta$, while the angle of $CBD$ at $C$ is $\pi-\gamma$. Applying the ideal-vertex case twice and subtracting gives
\[
\operatorname{area}(ABC)
=[\pi-\alpha-\beta-\delta]-[\pi-(\pi-\gamma)-\delta]
=\pi-\alpha-\beta-\gamma.
\]
\end{proof}

\begin{corollary}\label{geometry:cor:anglesum}
A hyperbolic triangle has angle sum strictly less than $\pi$ and area at most $\pi$. Its area equals $\pi$ exactly when all three vertices are ideal; otherwise its area is strictly less than $\pi$.
\end{corollary}

\begin{proof}
The nondegenerate triangle has positive area. Its angles are nonnegative, and an angle at a finite vertex is positive. Apply the area formula.
\end{proof}

The missing angle is an area: $\pi-(\alpha+\beta+\gamma)$ measures how much hyperbolic plane the triangle encloses. For instance, angles $\pi/2,\pi/3,\pi/7$ give area $\pi/42$. On a sphere the sign is reversed: area produces an angle excess. This already rules out a dilation that multiplies all lengths by a factor other than $1$: such a map would keep triangle angles, hence areas, while multiplying areas by the square of that factor. Curvature has supplied a length scale.

\section{The Poincar\'e disk}

The half-plane makes calculations at $\infty$ easy. For a picture with no preferred boundary point, use the \emph{Cayley map}
\[
C(z)=\frac{z-i}{z+i},\qquad C^{-1}(w)=i\frac{1+w}{1-w}.
\]
The inequality $|z-i|<|z+i|$ says exactly that $\operatorname{Im}z>0$, so $C$ maps $\mathbb H$ bijectively to $\mathbb D=\{|w|<1\}$. It puts $i$ at $0$.

\begin{proposition}[the disk model]\label{geometry:prop:disk}
In disk coordinates the length element is $2|dw|/(1-|w|^2)$, and
\[
\cosh d(w_1,w_2)=1+\frac{2|w_1-w_2|^2}{(1-|w_1|^2)(1-|w_2|^2)}.
\]
The h-lines through $0$ are the diameters, and $d(0,w)=2\operatorname{artanh}|w|$.
\end{proposition}

\begin{proof}
For $z=x+iy$,
\[
1-|C(z)|^2=\frac{4y}{|z+i|^2},\qquad |C'(z)|=\frac2{|z+i|^2}.
\]
Their quotient gives the length element. Also
$C(z_1)-C(z_2)=2i(z_1-z_2)/((z_1+i)(z_2+i))$. Substitution into the claimed distance expression leaves $|z_1-z_2|^2/(2y_1y_2)$, so the half-plane formula proves it. With $w_1=0$ it reads $\cosh d=(1+|w|^2)/(1-|w|^2)$, equivalent to $|w|=\tanh(d/2)$.

The map $C$ sends the imaginary axis to the real diameter, since $C(iy)=(y-1)/(y+1)$. Rotations $w\mapsto e^{i\vartheta}w$ preserve the disk length element and rotate this diameter to all the others. Their segments minimize length. Conversely a shortest arc through $0$ pulls back to the unique minimizing h-line arc through $i$; the preceding diameter construction supplies that arc for each other endpoint. Thus no additional h-lines through $0$ occur.
\end{proof}

A Euclidean radius approaching $1$ corresponds to hyperbolic distance tending to infinity. The boundary circle is outside the metric space, just as the real axis is outside $\mathbb H$. Both models preserve angles, so their curved-looking h-lines meet at the angles we see on the page.

\section{The isometry group}

Call a map $f:\mathbb H\to\mathbb H$ an \emph{isometry} if it preserves every distance. We do not assume it is onto. In this particular space, preservation of distances will force that conclusion and give a formula for the map.

\begin{lemma}\label{lem:hypcircles}
For $b>0$ and $r\ge0$, the hyperbolic circle $\{z:d(z,ib)=r\}$ is the Euclidean circle with centre $ib\cosh r$ and radius $b\sinh r$, understood as a point when $r=0$. Two distinct such circles meet in at most two points, symmetric about the imaginary axis.
\end{lemma}

\begin{proof}
The distance formula gives
\[
x^2+y^2-2by\cosh r+b^2=0,
\qquad x^2+(y-b\cosh r)^2=b^2\sinh^2r.
\]
The centre height $Y=b\cosh r$ and radius $R=b\sinh r$ determine $b=\sqrt{Y^2-R^2}$ and then $r\ge0$. Thus distinct pairs $(b,r)$ give distinct Euclidean circles. Subtracting their first equations either fixes $y$ or gives a nonzero constant equal to zero. A horizontal line cuts a circle in at most two points, of the form $\pm x+iy$. This includes tangency and radius-zero circles.
\end{proof}

\begin{theorem}[the isometry group]\label{thm:isometries}
Every isometry of $\mathbb H$ has the form $z\mapsto g(z)$ or $z\mapsto g(-\bar z)$ with $g\in PSL_2(\R)$. Consequently the orientation-preserving isometry group is $PSL_2(\R)$.
\end{theorem}

\begin{proof}
Normalize two image points first. The h-line lemma supplies $g$ sending $f(i),f(2i)$ to $i,it$ with $t>1$. Their distance is $\log2$, so $t=2$. Set $h=g\circ f$; it fixes $i,2i$.

For any $z$, the two distances to these fixed centres determine at most the pair $z,-\bar z$, by the circle lemma. The two circles are distinct because their hyperbolic centres differ. Hence $h(z)$ is one of these two points.

Suppose some off-axis point $z_0=x_0+iy_0$ is fixed. For $w=u+iv$ off the axis,
\[
|w-z_0|^2-|-\bar w-z_0|^2=-4ux_0\ne0.
\]
The denominators in the distance formula are the same for the two candidates. Only $w$ therefore has the required distance from $z_0$, and $h(w)=w$. Points on the axis are fixed in either case. Thus $h$ is the identity. If there is no fixed off-axis point, it is the reflection everywhere. Undoing the normalization proves the formulas and, in particular, surjectivity. The fractional linear maps preserve orientation, while reflection reverses it.
\end{proof}

\begin{proposition}[geodesic polar coordinates]\label{prop:hyppolar}
For $r>0$ and $\vartheta\in\R$, set $P(r,\vartheta)=C^{-1}(\tanh(r/2)e^{i\vartheta})$. Then $r=d(i,P(r,\vartheta))$, the curves of constant $\vartheta$ are the unit-speed geodesics from $i$, and the metric is
\[
\dd r^2+\sinh^2r\,\dd\vartheta^2.
\]
A circle of radius $r$ has length $2\pi\sinh r$; the disk it bounds has area $2\pi(\cosh r-1)$.
\end{proposition}

\begin{proof}
Work in the disk and put $\rho=\tanh(r/2)$. The distance formula gives $d(0,\rho e^{i\vartheta})=r$. Along a fixed diameter, increasing $r$ therefore runs at unit speed. Since $|dw|^2=d\rho^2+\rho^2d\vartheta^2$,
\[
\frac{4|dw|^2}{(1-\rho^2)^2}
=dr^2+4\sinh^2(r/2)\cosh^2(r/2)\,d\vartheta^2
=dr^2+\sinh^2r\,d\vartheta^2.
\]
At fixed $r$ integrate $\sinh r\,d\vartheta$ around $0\le\vartheta\le2\pi$. For area integrate $\sinh s\,ds\,d\vartheta$ over $0\le s\le r$.
\end{proof}

Angles are taken modulo $2\pi$. At $r=0$ every angle gives the same point, so these coordinates are not a chart at the centre; the length and area formulas still extend there with value $0$. Exercise~\ref{ex:geodlength} now has its promised example: $G=\sinh^2r$ gives $K=-(\sqrt G)''/\sqrt G=-1$. At large $r$, the circumference is about $\pi e^r$. The disk has exponentially more room than a Euclidean disk of the same radius.

\section{Points, directions and matrices}

A point takes two real coordinates; a unit direction at that point takes one more. The three parameters of $SL_2(\R)$ will record exactly this information. Use the subgroups
\[
K=SO(2),\qquad
A=\{a_t=\operatorname{diag}(e^{t/2},e^{-t/2}):t\in\R\},\qquad
N=\{n_x=\bigl(\begin{smallmatrix}1&x\\0&1\end{smallmatrix}\bigr):x\in\R\},
\]
and write $k_\theta=\bigl(\begin{smallmatrix}\cos\theta&-\sin\theta\\\sin\theta&\cos\theta\end{smallmatrix}\bigr)$. These are the rotation, diagonal and unipotent subgroups met in Chapter~7. Here $n_x$ moves a point horizontally and $a_t$ multiplies its height by $e^t$.

\begin{theorem}[Iwasawa decomposition]\label{thm:iwasawa}
\begin{enumerate}
\item[(a)] $SL_2(\R)$ acts transitively on $\mathbb H$, and the stabilizer of $i$ is $K$.
\item[(b)] Every $g\in SL_2(\R)$ has a unique expression $g=n_xa_tk_\theta$, with $x,t\in\R$ and $\theta\in\R/2\pi\Z$. These parameters depend smoothly on $g$.
\item[(c)] At $i$ the differential of $k_\theta$ is multiplication by $e^{-2i\theta}$; it rotates tangent vectors by $-2\theta$.
\end{enumerate}
\end{theorem}

\begin{proof}
The map $n_xa_t$ sends $i$ to $x+ie^t$, covering every point of $\mathbb H$. A matrix fixes $i$ precisely when $ai+b=i(ci+d)$, or $a=d$, $b=-c$; its determinant condition then says $a^2+b^2=1$. These are exactly the matrices of $K$, proving (a).

For (b), write $g(i)=x+iy$ and take $t=\log y$. The matrix $(n_xa_t)^{-1}g$ fixes $i$ and hence equals $k_\theta$. Conversely the image of $i$ determines $x,t$, after which the remaining matrix determines $\theta$ modulo $2\pi$. For $g=\bigl(\begin{smallmatrix}a&b\\c&d\end{smallmatrix}\bigr)$ the formulas are
\[
x=\frac{ac+bd}{c^2+d^2},\qquad e^t=\frac1{c^2+d^2},\qquad k_\theta=a_{-t}n_{-x}g.
\]
They are smooth, with the last parameter regarded as a point of the circle rather than a globally chosen real angle. Finally $k_\theta'(i)=(i\sin\theta+\cos\theta)^{-2}=e^{-2i\theta}$, giving (c).
\end{proof}

The factor $2$ and its minus sign matter: our matrix parameter $\theta$ is not the physical angle of a tangent vector. In $PSL_2(\R)$, $\theta$ is taken modulo $\pi$, so the physical direction makes exactly one turn.

The \emph{unit tangent bundle} $T^1\mathbb H$ consists of pairs $(z,v)$ in which $v$ is a tangent vector of hyperbolic length $1$. Its Euclidean length is $\operatorname{Im}z$. Let $v_0=i$ be the upward unit vector at $i$. The \emph{geodesic flow} follows the unique unit-speed h-line with prescribed initial point and vector:
\[
\phi_t(z,v)=(\gamma(t),\gamma'(t)),\qquad \gamma(0)=z,\quad\gamma'(0)=v.
\]
Isometries preserve the lengths and uniqueness of minimizing arcs, so applying an isometry before following the geodesic gives the same result as applying it afterwards.

\begin{theorem}[the geodesic flow is a matrix product]\label{thm:geodflow}
The map $\Psi:PSL_2(\R)\to T^1\mathbb H$ given by $\Psi(g)=(g(i),dg_i(v_0))$ is a bijection, and
\[
\phi_t(\Psi(g))=\Psi(ga_t).
\]
\end{theorem}

\begin{proof}
The matrices taking $i$ to a chosen point form a coset $gK$. Within this coset the derivative of $k_\theta$ rotates the direction through all possibilities, exactly once as $\theta$ ranges modulo $\pi$. Thus a point and a unit vector determine exactly one element of $PSL_2(\R)$.

At the base point, $a_t(i)=ie^t$ is a vertical unit-speed geodesic, with velocity $ie^t=d(a_t)_i(v_0)$. Hence $\phi_t(\Psi(I))=\Psi(a_t)$. Applying $g$ to this identity and using the commutation with isometries gives $\phi_t(\Psi(g))=\Psi(ga_t)$.
\end{proof}

Right multiplication by $n_s$ gives a different flow. At the identity its base point is $n_s(i)=s+i$, a horizontal curve of constant height. Its carried vector is upward, perpendicular to the curve, rather than tangent to it. An isometric image of such a horizontal line is called a \emph{horocycle}; the marked vector is its chosen unit normal. Thus right multiplication by $N$ is the horocycle flow on the space of marked points and directions. Left multiplication moves the whole geometric picture, while right multiplication by $A$ or $N$ follows one of these two motions within it.

A discrete subgroup $\Gamma$ acts by left multiplication, and both right flows descend to $\Gamma\backslash PSL_2(\R)$. If the action on $\mathbb H$ has no fixed points, this is the unit tangent bundle of the quotient surface. With elliptic fixed points, as for the modular group below, the quotient has cone points and is an \emph{orbifold}: locally a disk divided by a finite rotation group. Its unit tangent space is still $\Gamma\backslash PSL_2(\R)$; the left action on the group is free. We do not treat the cone points as smooth surface points.

Hopf's ergodicity theorem, which we quote, describes the geodesic flow for finite-area hyperbolic quotients such as the modular example. Almost every initial point and direction spends in a measurable region the fraction of time specified by its normalized invariant volume. Volume~11 develops the ergodic theory needed to prove such assertions; the matrix description here is their geometric starting point.

\begin{theorem}[elliptic, parabolic, hyperbolic]\label{thm:trace}
Let $g\in SL_2(\R)$ with $g\ne\pm I$.
\begin{enumerate}
\item[(a)] If $|\operatorname{tr}g|<2$, then $g$ has exactly one fixed point $z_0\in\mathbb H$ and is conjugate to some $k_\theta$. It is a rotation about $z_0$, called \emph{elliptic}.
\item[(b)] If $|\operatorname{tr}g|=2$, then $g$ has no fixed point in $\mathbb H$ and exactly one in $\R\cup\{\infty\}$. It is conjugate to $\pm n_1$ or $\pm n_{-1}$ and is called \emph{parabolic}.
\item[(c)] If $|\operatorname{tr}g|>2$, then $g$ is conjugate to $\pm\operatorname{diag}(\lambda,\lambda^{-1})$, where $\lambda>1$, and is called \emph{hyperbolic}. It fixes two boundary points and translates their joining h-line, its \emph{axis}, by $\ell=2\log\lambda$. For every $z\in\mathbb H$,
\[
d(z,g(z))\ge\ell,
\qquad 2\cosh(\ell/2)=|\operatorname{tr}g|,
\]
with equality in the distance inequality exactly on the axis.
\end{enumerate}
\end{theorem}

\begin{proof}
A finite fixed point satisfies $cz^2+(d-a)z-b=0$. Its discriminant is $(a+d)^2-4$, so the trace divides the argument into three cases.

If its absolute value is less than $2$, then $c\ne0$: otherwise $ad=1$ would give $|a+d|\ge2$. The two roots are conjugate and nonreal, exactly one in $\mathbb H$. Move that root to $i$ by conjugation. The conjugated matrix fixes $i$, so the Iwasawa theorem puts it in $K$.

If the absolute trace exceeds $2$, the matrix has two distinct real eigenvalues. After replacing $g$ by $-g$ if necessary, order them as $\lambda,\lambda^{-1}$ with $\lambda>1$. Choose eigenvectors in this order, change the sign of one if needed, and scale the resulting basis matrix to have determinant $1$. In that basis the transformation is $z\mapsto\lambda^2z$. It fixes $0,\infty$ and moves the imaginary axis a distance $2\log\lambda$. For $z=x+iy$,
\[
\cosh d(z,\lambda^2z)
=1+\frac{(\lambda^2-1)^2(x^2+y^2)}{2\lambda^2y^2}
\ge\frac{\lambda^2+\lambda^{-2}}2
=\cosh(2\log\lambda).
\]
Equality holds exactly when $x=0$. Conjugation preserves distances and sends this axis to the axis of $g$. The trace identity follows from $\lambda+\lambda^{-1}=2\cosh(\log\lambda)$.

Finally, absolute trace $2$ gives a repeated eigenvalue $\varepsilon\in\{1,-1\}$. Since $g\ne\varepsilon I$, choose a real eigenvector $v$ and complete it to an oriented determinant-one basis $(v,w)$. In this basis $g=\varepsilon\bigl(\begin{smallmatrix}1&u\\0&1\end{smallmatrix}\bigr)$ with $u\ne0$. Conjugating by $a_t$ replaces $u$ by $e^tu$, so its magnitude can be made $1$. Translation by $1$ or $-1$ fixes only $\infty$ on the boundary and no point in $\mathbb H$.
\end{proof}

\section{The modular group}

Now restrict the matrix entries to integers. Volume~6, Proposition~12.1 proved that $SL_2(\Z)$ is generated by
\[
T=\begin{pmatrix}1&1\\0&1\end{pmatrix},\qquad
S=\begin{pmatrix}0&-1\\1&0\end{pmatrix}.
\]
They act as $T(z)=z+1$ and $S(z)=-1/z$. Volume~6, Theorem~12.2 puts every orbit in
\[
F=\{z\in\mathbb H:|\operatorname{Re}z|\le\tfrac12,\ |z|\ge1\}.
\]
Distinct interior points of $F$ are never equivalent. Equivalent boundary points are paired by $T,T^{-1}$ or $S$. Such a region is called a \emph{fundamental domain}: its translates cover the plane, overlapping only along boundaries. Figure~\ref{fig:modular} shows the resulting tiles.

The vertices of $F$ are $\rho=e^{2\pi i/3}$, $\rho+1=e^{i\pi/3}$ and the ideal point $\infty$. At each finite vertex the angle between the vertical side and the unit circle is $\pi/3$. Therefore
\[
\operatorname{area}(F)=\pi-\pi/3-\pi/3=\pi/3.
\]
Identify its two vertical sides by $T$ and the two halves of its circular side by $S$. The quotient is called the \emph{modular surface}, though its cone points make it an orbifold. Its underlying shape is a sphere with one cusp, an end extending to infinite height. The cone angles are $\pi$ at $i$ and $2\pi/3$ at $\rho$. Its area is finite despite the infinite cusp.

A hyperbolic integer matrix identifies points along its axis. The projected axis is therefore a closed geodesic traversed in length $\ell$, with $2\cosh(\ell/2)=|\operatorname{tr}g|$. If the matrix is a power of a smaller translation in the group, this traversal goes around the underlying closed geodesic more than once. Since its absolute trace is an integer greater than $2$, Exercise~\ref{ex:shortest} gives a lower bound for these lengths.

\begin{figure}[ht]
\centering
\fig{ch09-modular}
\caption{The tessellation of $\mathbb H$ by the images of the fundamental domain $F$ (shaded) under $SL_2(\Z)$. The tiles accumulate at every rational point of the real axis; each is a
hyperbolic triangle of area $\frac\pi3$, and they are all congruent, however small some of them look.}
\label{fig:modular}
\end{figure}

\section{Farey tessellation and cutting sequences}

Write a rational boundary point as $p/q$ in lowest terms with $q>0$, and write $\infty=1/0$. Two points $p/q,r/s$ are \emph{Farey neighbours} when $|ps-qr|=1$. Their joining h-line is a \emph{Farey edge}. For example, the edges ending at $\infty$ are precisely the vertical lines through integers: the determinant condition forces the other denominator to be $1$.

The first open triangle to keep in mind has vertices $0,1,\infty$:
\[
\Delta_0=\{z\in\mathbb H:0<\operatorname{Re}z<1,\ |z-\tfrac12|>\tfrac12\}.
\]
Its three sides are Farey edges. The determinant condition prevents other Farey edges from running through its interior.

\begin{proposition}[the Farey tessellation]\label{prop:farey}
\begin{enumerate}
\item[(a)] $GL_2(\Z)$ permutes Farey edges. A matrix acts by $(az+b)/(cz+d)$ when its determinant is $1$, and by $(a\bar z+b)/(c\bar z+d)$ when its determinant is $-1$. These are isometries of $\mathbb H$, and $SL_2(\Z)$ acts transitively on Farey edges.
\item[(b)] Distinct Farey edges do not cross.
\item[(c)] No Farey edge meets $\Delta_0$ or any $g(\Delta_0)$ with $g\in SL_2(\Z)$. If $ps-qr=-1$, the triangles on the two sides of the edge from $p/q$ to $r/s$ have third vertices $(p+r)/(q+s)$ and $(p-r)/(q-s)$, respectively. A zero denominator denotes $\infty$.
\end{enumerate}
\end{proposition}

\begin{proof}
For (a), a matrix with integral inverse takes primitive integer column vectors to primitive vectors. Applying it to the two endpoint columns multiplies their determinant by $\pm1$, preserving the neighbour condition. A determinant-minus-one transformation is reflection $z\mapsto-\bar z$ followed by the determinant-one matrix obtained by negating its first column. It is therefore an isometry. If $ps-qr=1$, the matrix with columns $(p,q)^T,(r,s)^T$ sends $\infty,0$ to the chosen endpoints. Interchange the endpoint columns when the determinant is $-1$. This proves transitivity for unoriented edges.

For (b), two distinct h-lines cross exactly when their endpoints alternate around the boundary circle. This follows directly for a vertical axis, whose opposite sides are the positive and negative real half-lines, and then for every h-line by an isometry. Send one proposed crossing edge to $(0,\infty)$. The other would have finite endpoints $p/q<0<r/s$, with $p<0<r$ and $q,s\ge1$. But then
\[
|ps-qr|=|p|s+qr\ge2,
\]
contradicting the Farey condition.

For (c), a distinct edge entering $\Delta_0$ could not cross any of its three sides, by (b). It would have to remain there in both directions. The closure of this triangle meets the boundary only at $0,1,\infty$, so the edge would have to join two of these vertices and hence be a side, a contradiction. Part (a) transports this argument to every image. Finally
\[
g=\begin{pmatrix}r&p\\s&q\end{pmatrix}\in SL_2(\Z)
\]
sends $0,\infty$ to $p/q,r/s$. The triangles $\Delta_0$ and $\Delta_0-1$ share $(0,\infty)$ and have third vertices $1,-1$. Their images have the claimed third vertices, since $g(1)=(p+r)/(q+s)$ and $g(-1)=(r-p)/(s-q)$.
\end{proof}

Each Farey triangle is ideal and has area $\pi$, three times the area of $F$. The Farey edges give a coarser pattern than the modular tiles. For the coding argument we need just the triangles displayed in part (c) and the fact that no further edge cuts through them.

Orient an h-line $\gamma$ from a finite endpoint $y$ to a finite endpoint $x$. The boundary has two arcs between these endpoints. Its \emph{right arc} is $(y,x)$ when $y<x$ and the arc through $\infty$ when $y>x$. This agrees with the right-hand side as you travel along $\gamma$. The other arc is the left arc. Whenever $\gamma$ enters a Farey triangle through one edge and leaves through another, those edges share a vertex. Record $L$ if that vertex lies on the left arc and $R$ if it lies on the right. These letters form the \emph{cutting sequence}.

\begin{theorem}[Series, 1985]\label{thm:cutting}
Let $x>1$ be irrational, with continued fraction $x=[a_0;a_1,a_2,\ldots]$, and let $-1<y<0$. After crossing the edge $(0,\infty)$, the h-line from $y$ to $x$ has cutting sequence
\[
L^{a_0}R^{a_1}L^{a_2}R^{a_3}\cdots.
\]
\end{theorem}

\begin{proof}
Put $m=a_0=\lfloor x\rfloor$. Since $y<k<x$ for $0\le k\le m$, the path crosses each vertical edge $(k,\infty)$. Between consecutive ones it crosses $\Delta_0+k$, with no intervening edge by the Farey proposition. The shared vertex is $\infty$, on the left. These first $m$ triangles give $L^m$.

After crossing $(m,\infty)$ it enters $\Delta_0+m$ and leaves through $(m,m+1)$, because $m<x<m+1$. That triangle's shared vertex is $m$, on the right, so the block of $L$'s has ended. To count the next block, apply
\[
\psi(z)=\frac1{\bar z-m}.
\]
This determinant-minus-one integral transformation is an isometry preserving Farey edges. It sends $m,\infty,m+1$ to $\infty,0,1$. The transformed path goes from
\[
y'=\frac1{y-m}\in(-1,0)\quad\hbox{to}\quad
x'=\frac1{x-m}=[a_1;a_2,\ldots]>1.
\]
The same starting configuration has returned, but reflection reverses left and right. More explicitly, $\psi$ takes $(y,m)$ to $(-\infty,y')$ and $(m,x)$ to $(x',\infty)$; together these are the left arc of the new path. Thus its initial $L^{a_1}$ corresponds to $R^{a_1}$ in the original picture. Repeat at each new continued-fraction digit. Every finite initial block is obtained this way, proving the infinite sequence.
\end{proof}

\begin{figure}[ht]
\centering
\fig{ch09-farey}
\caption{The Farey tessellation (edges between Farey neighbours, here with denominators up to $12$) and the h-line from $1-\sqrt2$ to $1+\sqrt2=[2;2,2,\dots]$. After crossing the edge from $0$ to $\infty$ it
passes two triangles with the vertex $\infty$ on its left, then two with a vertex on its right, and so on: the cutting sequence $LLRRLLRR\cdots$.}
\label{fig:farey}
\end{figure}

\begin{figure}[ht]
\centering\fig{v9-ch09-geodesic-horocycle}
\caption{At $i$, the geodesic flow moves upward along $ie^t$. The horocycle flow moves horizontally along $s+i$ while carrying the upward unit normal. The carried vector is not the horocycle's tangent.}
\label{fig:v9-ch09-geodesic-horocycle}
\end{figure}

\begin{figure}[ht]
\centering\fig{v9-ch09-axis-displacement}
\caption{For $z\mapsto e^\ell z$, the imaginary axis attains the minimum displacement $\ell$. At the same height, moving away from the axis increases $\cosh d$ by $(\cosh\ell-1)(x/y)^2$. The drawn example uses $\ell=\log2$.}
\label{fig:v9-ch09-axis-displacement}
\end{figure}

\begin{beauty}
The operation $x\mapsto1/(x-\lfloor x\rfloor)$ from continued fractions is now an isometry that resets a path in a triangle. The digits count successive runs through the tiling. A hyperbolic integer matrix fixes two real quadratic conjugates: its fixed-point equation has integral coefficients and two real roots. The h-line joining that pair projects to a closed geodesic. One cannot pair a quadratic irrational with an arbitrary second endpoint and expect the path to close. Exercises~\ref{ex:shortest} and~\ref{ex:sqrt2axis} show two explicit conjugate pairs and their periods.
\end{beauty}

\begin{remark}
Artin used the connection between continued fractions and geodesics in 1924 to construct a dense geodesic trajectory on the modular surface. The coding and ergodic results beyond the theorem above are quoted here. Series' work describes the geodesic flow through a suitable return map and a two-sided continued-fraction coding; boundary identifications and exceptional trajectories need care. The Gauss map keeps only the forward information. Its invariant probability density $dx/((1+x)\log2)$ is related to Haar measure on the space of points and directions. Volume~11 returns to this measure-theoretic connection.
\end{remark}

\section{What we proved, and what comes next}

Try reading the example $1+\sqrt2=[2;2,2,\ldots]$ from Figure~\ref{fig:farey}: two triangles share a vertex to the left of the path, then two share one to the right. The repeated digits become the repeated block $LLRR$. The relevant matrix, found in Exercise~\ref{ex:sqrt2axis}, translates that same path by one period. A continued fraction, a closed geodesic and an integer matrix are three descriptions of this example.

To follow the construction again, start with the distance formula and the three trace cases. They explain what the lines and motions mean. Then use the determinant condition $|ps-qr|=1$ to find the triangles; the cutting proof needs only this condition and the continued-fraction step. The area computation is a separate useful tool, especially for the exercises about the cusp and Klein's quartic.

All the orientation-preserving symmetries used here are holomorphic maps. The next chapter asks how far holomorphic maps can reshape a region while preserving angles locally. The Riemann mapping theorem answers for nonempty simply connected proper regions of the plane: each is conformally equivalent to the disk.

\section{Exercises}

\begin{exercise}[$\star$]\label{ex:iwasawaexplicit}
(a) Factor $\bigl(\begin{smallmatrix}2&1\\1&1\end{smallmatrix}\bigr)$ as $n_xa_tk_\theta$. (b) Show that the curves $t\mapsto k_\theta a_tk_\theta^{-1}(i)$ give all geodesics through $i$, and that $s\mapsto k_\theta n_sk_\theta^{-1}(i)$ gives all horocycles through $i$. Describe both families after applying the Cayley map.
\end{exercise}

\begin{exercise}[$\star\star$]\label{ex:shortest}
(a) Prove that a hyperbolic $g\in SL_2(\Z)$ satisfies $|\operatorname{tr}g|\ge3$ and has translation length
$\ell\ge2\operatorname{arccosh}(3/2)=4\log\varphi$, where $\varphi=(1+\sqrt5)/2$.
(b) Show that $\bigl(\begin{smallmatrix}2&1\\1&1\end{smallmatrix}\bigr)$ attains the bound, and find its axis.
(c) Orient the axis towards $\varphi$ and find its cutting sequence.
\end{exercise}

\begin{exercise}[$\star$]\label{ex:sqrt2axis}
Find a matrix in $SL_2(\Z)$ whose axis joins $1-\sqrt2$ to $1+\sqrt2$, as in Figure~\ref{fig:farey}. Compute its translation length and verify that it shifts the crossed triangles by one full block $LLRR$.
\end{exercise}

\begin{exercise}[$\star\star$]\label{ex:formsreduction}
Let $Q(x,y)=ax^2+bxy+cy^2$ be positive definite, with $a>0$ and $D=b^2-4ac<0$, and set $\tau_Q=(-b+\sqrt D)/(2a)\in\mathbb H$.
(a) For $g\in SL_2(\Z)$, prove that the root attached to $Q\circ g$ is $g^{-1}(\tau_Q)$.
(b) Prove that $\tau_Q\in F$ exactly when $|b|\le a\le c$; such a form is called \emph{reduced}.
(c) For integral forms, list every reduced triple $(a,b,c)$ of discriminant $-20$ and $-23$, then identify boundary duplicates. The numbers of classes are $2$ and $3$.
\end{exercise}

\begin{exercise}[$\star$]\label{ex:klein336}
(a) Find the area of a hyperbolic triangle with angles $\pi/2,\pi/3,\pi/7$.
(b) Klein's quartic is a compact genus-$3$ surface of curvature $-1$ and has a tiling by such triangles. Use Corollary~\ref{cor:gaussbonnet} to count them. For context, the quoted Hurwitz bound of 1893 gives at most $84(g-1)$ conformal automorphisms in genus $g\ge2$, and the quartic attains it. Distinguish the area calculation from these quoted facts about the tiling and its symmetries.
\end{exercise}

\begin{exercise}[$\star\star$]\label{ex:horocyclearea}
Put $B_h=\{z\in\mathbb H:\operatorname{Im}z>h\}$, called a \emph{horoball}.
(a) Compute $\operatorname{area}(F\cap B_h)$ for $h\ge1$.
(b) Prove $d(i+s,i)=2\operatorname{arsinh}(|s|/2)$ and compare this distance with the length $|s|$ along the horizontal horocycle. On a circle of radius $r$, compare the distance between antipodal points with a semicircular arc's length.
\end{exercise}

\begin{exercise}[$\star\star$]\label{ex:v9-ch09-axis-distance}
Let $g(z)=e^\ell z$, where $\ell>0$, and let $\rho$ be the distance from $z=x+iy$ to its axis.
(a) Find the unique closest axis point and prove $\sinh\rho=|x|/y$.
(b) Deduce
\[
\cosh d(z,g(z))=\cosh\ell+(\cosh\ell-1)\sinh^2\rho.
\]
Explain both the equality case on the axis and the growth of displacement far from it.
\end{exercise}

\begin{exercise}[$\star\star$]\label{ex:v9-ch09-ford-horoballs}
Fix $h\ge1$.
(a) If $g=\bigl(\begin{smallmatrix}a&b\\c&d\end{smallmatrix}\bigr)\in SL_2(\Z)$ with $c\ne0$, show that $g(B_h)$ is the open Euclidean disk tangent to the real axis at $a/c$, with radius $1/(2hc^2)$.
(b) Prove that the disks for distinct reduced fractions have disjoint interiors; equality of the tangency test at $h=1$ occurs exactly for Farey neighbours. Show also that their interiors miss $B_h$.
(c) Find the area removed from the modular quotient by truncating its cusp at height $h$, and the area remaining.
\end{exercise}

\begin{exercise}[$\star\star$]\label{ex:v9-ch09-axis-centralizer}
Fix $a_\ell\in PSL_2(\R)$ with $\ell>0$.
(a) Find all orientation-preserving isometries commuting with it.
(b) Find all $h\in PSL_2(\R)$ such that $ha_\ell h^{-1}=a_{-\ell}$.
(c) Transport both answers by conjugation to an arbitrary hyperbolic element. Explain the difference between preserving its two axis endpoints individually and exchanging them.
\end{exercise}

\begin{exercise}[$\star\star\star$]\label{ex:v9-ch09-direction-measure}
Use coordinates $(x,y,\alpha)$ on $T^1\mathbb H$, where the vector is $y e^{i\alpha}$ and $\alpha$ is modulo $2\pi$.
(a) Show that $dx\,dy\,d\alpha/y^2$ is invariant under the left action of $PSL_2(\R)$, and hence is Haar volume under $\Psi$.
(b) Prove that this group is unimodular by calculating the determinant of its adjoint action on $\mathfrak{sl}_2(\R)$. You may use the determinant criterion from Chapter~8. Deduce invariance under both right flows.
(c) Find the total volume of the modular unit tangent quotient and the probability of a base point lying above height $h\ge1$ in $F$, with direction restricted to an arc of length $\beta\in[0,2\pi]$. This computes a probability, not an ergodicity proof.
\end{exercise}

\section{Hints}

\begin{hint}{ex:iwasawaexplicit}
Compute $g(i)$ first. Then remove $n_xa_t$ on the left. In the disk, rotate the real diameter and the circle $|w-1/2|=1/2$.
\end{hint}

\begin{hint}{ex:shortest}
Use the integer trace and the trace--length identity. The fixed-point equation is $z^2-z-1=0$; the positive root has all continued-fraction digits equal to $1$.
\end{hint}

\begin{hint}{ex:sqrt2axis}
Require $cz^2+(d-a)z-b$ to be a multiple of $z^2-2z-1$. Try $c=b=2$. Factor the resulting matrix as the square of a determinant-minus-one integer matrix.
\end{hint}

\begin{hint}{ex:formsreduction}
Factor $Q(x,y)=a(x-\tau_Qy)(x-\bar\tau_Qy)$. For enumeration, $3a^2\le4ac-b^2=|D|$, and $b$ has the parity of $D$.
\end{hint}

\begin{hint}{ex:klein336}
Subtract the three angles from $\pi$, then divide the surface area by the tile area. Area alone does not construct the tiling or the automorphisms.
\end{hint}

\begin{hint}{ex:horocyclearea}
For $h\ge1$ the cusp is a strip of width $1$. For distances, use $\cosh(2u)=1+2\sinh^2u$. A diameter through the centre joins antipodal points by a minimizing arc.
\end{hint}

\begin{hint}{ex:v9-ch09-axis-distance}
Minimize $\cosh d(z,it)=(x^2+y^2+t^2)/(2yt)$ over $t>0$. Then substitute $(x/y)^2=\sinh^2\rho$ into the displacement formula.
\end{hint}

\begin{hint}{ex:v9-ch09-ford-horoballs}
For $w=u+iv$, use $\operatorname{Im}(g^{-1}w)=v/|-cw+a|^2$. For two disks, subtract the square of the sum of their radii from the squared distance between their centres.
\end{hint}

\begin{hint}{ex:v9-ch09-axis-centralizer}
Work with real determinant-one representatives. A possible sign in the commuting equation is ruled out by the nonzero trace. Compare entries against a diagonal matrix with distinct eigenvalues.
\end{hint}

\begin{hint}{ex:v9-ch09-direction-measure}
The new angle is $\alpha+\arg g'(z)$, so the full Jacobian has a triangular block form. For the adjoint action use $E=\bigl(\begin{smallmatrix}0&1\\0&0\end{smallmatrix}\bigr)$, $H=\operatorname{diag}(1,-1)$ and $F=\bigl(\begin{smallmatrix}0&0\\1&0\end{smallmatrix}\bigr)$. Iwasawa decomposition reduces the determinant check to $N,A,K$.
\end{hint}

\section{Solutions}

\begin{solution}{ex:iwasawaexplicit}
\idea Locate the image of $i$ before recovering the remaining rotation.
\writeup (a) The Iwasawa formulas give $x=3/2$, $e^t=1/2$. Removing these factors leaves
\[
a_{\log2}n_{-3/2}\begin{pmatrix}2&1\\1&1\end{pmatrix}
=\frac1{\sqrt2}\begin{pmatrix}1&-1\\1&1\end{pmatrix}=k_{\pi/4}.
\]
Thus $g=n_{3/2}a_{-\log2}k_{\pi/4}$, and indeed $g(i)=3/2+i/2$.
(b) Since $k_\theta$ fixes $i$, the curves are $k_\theta(ie^t)$ and $k_\theta(i+s)$. Its derivative rotates all directions, so these give all geodesics and all horocycles through $i$. The Cayley map takes the first family to the disk diameters. For the second, $\operatorname{Im}C^{-1}(w)=(1-|w|^2)/|1-w|^2=1$ is equivalent to $|w-1/2|=1/2$, omitting its boundary point $1$. Rotations give all such horocycles through $0$.
\end{solution}

\begin{solution}{ex:shortest}
\idea The trace is integral, so it cannot approach $2$ from above without a gap.
\writeup (a) Hyperbolicity gives an integer $|\operatorname{tr}g|>2$, hence at least $3$. The trace formula gives $\ell\ge2\operatorname{arccosh}(3/2)$. Solving $\cosh u=3/2$ gives $e^u=(3+\sqrt5)/2=\varphi^2$, so the bound is $4\log\varphi\approx1.9248$.
(b) The stated matrix has trace $3$. Its fixed points solve $z^2-z-1=0$, giving $\varphi$ and $1-\varphi=-1/\varphi$. The axis is their joining semicircle. The vector $(\varphi,1)^T$ has eigenvalue $\varphi^2>1$, so $\varphi$ is the attracting endpoint.
(c) Since $\varphi=[1;1,1,\ldots]$ and $-1/\varphi\in(-1,0)$, the cutting theorem gives $LRLR\cdots$. The determinant-minus-one matrix $M=\bigl(\begin{smallmatrix}1&1\\1&0\end{smallmatrix}\bigr)$ acts as $1+1/\bar z$ and undoes one continued-fraction reset. It shifts one block while reversing left and right. Its square is $g$, so $g$ shifts one full period $LR$.
\end{solution}

\begin{solution}{ex:sqrt2axis}
\idea Match the desired quadratic equation to the matrix's fixed-point equation.
\writeup Taking $c=b=k$ and $d=a-2k$ gives the required roots. The determinant equation becomes $(a-k)^2-2k^2=1$. Choose $k=2$, $a=5$, $d=1$:
\[
g=\begin{pmatrix}5&2\\2&1\end{pmatrix},\qquad
\lambda=3+2\sqrt2=(1+\sqrt2)^2,\qquad
\ell=4\log(1+\sqrt2)\approx3.5255.
\]
No minimality claim is needed to supply this example. The Pell equation is the one studied in Volume~2, Theorem~9.8. Now $g=M^2$ for $M=\bigl(\begin{smallmatrix}2&1\\1&0\end{smallmatrix}\bigr)$, acting as $2+1/\bar z$. It is the inverse reset for digit $2$, preserves the tessellation and fixes both endpoints of the axis. Each application advances a two-letter block and exchanges left and right; two applications advance $LLRR$. This agrees with $1+\sqrt2=[2;2,2,\ldots]$ and the cutting theorem.
\end{solution}

\begin{solution}{ex:formsreduction}
\idea Moving a form's root into $F$ translates reduction inequalities into geometric ones.
\writeup (a) Write $\tau=\tau_Q$ and $g=\bigl(\begin{smallmatrix}p&q\\r&s\end{smallmatrix}\bigr)$. The factorization gives a root of $Q(px+qy,rx+sy)$ at
\[
\frac{x}{y}=\frac{s\tau-q}{p-r\tau}=g^{-1}(\tau).
\]
The denominator is nonzero in $\mathbb H$, and the inverse transformation keeps the root in $\mathbb H$.
(b) We have $\operatorname{Re}\tau=-b/(2a)$ and $|\tau|^2=c/a$. Thus the two conditions defining $F$ are exactly $|b|\le a\le c$.
(c) These inequalities give $|D|=4ac-b^2\ge3a^2$. For $D=-20$, $a\le2$ and $b$ is even. The full list is
\[
(1,0,5),\quad(2,2,3),\quad(2,-2,3).
\]
The last two roots differ by $1$ and lie on the vertical boundary of $F$, so they form one class. The first is interior, giving two classes.
For $D=-23$, again $a\le2$, now with $b$ odd. The full list is
\[
(1,1,6),\quad(1,-1,6),\quad(2,1,3),\quad(2,-1,3).
\]
The first two roots are boundary translates and form one class. The last two are distinct interior points of $F$, hence give distinct classes. There are three classes. Listing all triples before making identifications avoids dropping one reduced representative.
\end{solution}

\begin{solution}{ex:klein336}
\idea Once a tiling is given, additivity of area counts its tiles.
\writeup (a) The area is $\pi(1-1/2-1/3-1/7)=\pi/42$.
(b) Gauss--Bonnet gives $-\operatorname{area}=2\pi(2-2\cdot3)=-8\pi$, so there are $8\pi/(\pi/42)=336$ triangles. In the quoted tiling of the quartic, the triangles have two orientation types, with $168$ of each. The known automorphism action permutes each type simply transitively and attains $84(3-1)=168$. This last statement uses the geometry of the quartic; the area division by itself proves only the count $336$.
\end{solution}

\begin{solution}{ex:horocyclearea}
\idea High in the cusp the region has constant Euclidean width, but the area density falls as $y^{-2}$.
\writeup (a) For $h\ge1$ the region is $-1/2\le x\le1/2$, $y>h$, so
\[
\operatorname{area}(F\cap B_h)=\int_{-1/2}^{1/2}\int_h^\infty y^{-2}\,dy\,dx=1/h.
\]
(b) The distance formula gives $\cosh d=1+s^2/2$. Since distances are nonnegative, the double-angle identity yields $d=2\operatorname{arsinh}(|s|/2)$. The horocyclic arc has length $|s|$, while this distance is asymptotic to $2\log|s|$ as $|s|\to\infty$. For antipodal points on a radius-$r$ circle, the diameter through the centre is a minimizing geodesic of length $2r$. Either semicircular arc has length $\pi\sinh r$, asymptotic to $(\pi/2)e^r$. At $r=0$ both lengths are zero.
\end{solution}

\begin{solution}{ex:v9-ch09-axis-distance}
\idea The closest point is found by minimizing a one-variable expression, not by assuming a Euclidean perpendicular.
\writeup (a) Put $R=\sqrt{x^2+y^2}>0$. The distance formula gives
\[
\cosh d(z,it)=\frac{R^2+t^2}{2yt}\ge\frac Ry,
\]
with equality exactly at $t=R$, since $R^2+t^2-2Rt=(t-R)^2$. Thus the unique closest point is $iR$ and $\cosh\rho=R/y$. Consequently $\sinh^2\rho=R^2/y^2-1=x^2/y^2$, giving the nonnegative square root $\sinh\rho=|x|/y$.
(b) Direct substitution gives
\[
\cosh d(z,e^\ell z)
=1+\frac{(e^\ell-1)^2}{2e^\ell}\left(1+\frac{x^2}{y^2}\right)
=\cosh\ell+(\cosh\ell-1)\sinh^2\rho.
\]
The extra term vanishes exactly when $x=0$, or $\rho=0$. For large $\rho$, it is asymptotic to $(\cosh\ell-1)e^{2\rho}/4$, so displacement tends to infinity and is $2\rho+O(1)$ for fixed $\ell>0$.
\end{solution}

\begin{solution}{ex:v9-ch09-ford-horoballs}
\idea Images of a cusp are ordinary tangent disks, so an integer determinant controls whether they overlap.
\writeup (a) Write $w=u+iv$. Membership in $g(B_h)$ is equivalent to
\[
\frac{v}{|-cw+a|^2}>h,
\]
which, for $c\ne0$, becomes
\[
(u-a/c)^2+\left(v-\frac1{2hc^2}\right)^2
<\frac1{4h^2c^4}.
\]
This disk is in $\mathbb H$ and tangent at $a/c$. Fractions may be represented with positive denominators since changing both signs changes nothing.
(b) For distinct reduced fractions $a/c,a'/c'$ with $c,c'>0$, put $R=1/(2hc^2)$ and $R'=1/(2hc'^2)$. The centres are $(a/c,R)$ and $(a'/c',R')$. Their squared distance minus $(R+R')^2$ is
\[
\left(\frac ac-\frac{a'}{c'}\right)^2-4RR'
=\frac{(ac'-a'c)^2-h^{-2}}{c^2c'^2}\ge0.
\]
The integer numerator determinant is nonzero, hence has absolute value at least $1$. This proves disjoint interiors; at $h=1$, equality means exactly Farey neighbours. Each disk has top height $2R=1/(hc^2)\le h$, so its interior misses $B_h$ as well.
(c) Boundary identifications have area zero. The removed cusp area is $1/h$ and the remaining quotient area is $\pi/3-1/h$. This is positive for $h\ge1$ and approaches $\pi/3$ as $h\to\infty$.
\end{solution}

\begin{solution}{ex:v9-ch09-axis-centralizer}
\idea Distinct eigenvalues force a commuting map to preserve the two eigendirections.
\writeup Let $D=\operatorname{diag}(e^{\ell/2},e^{-\ell/2})$.
(a) Commutation in $PSL_2(\R)$ says $HDH^{-1}=\pm D$ for determinant-one representatives. Taking traces excludes the minus sign. Comparing entries in $HD=DH$ forces the off-diagonal entries of $H$ to vanish. Thus $H=\operatorname{diag}(r,r^{-1})$, $r\ne0$. Modulo its negative, take $r>0$ and obtain precisely $a_t$, $t=2\log r\in\R$.
(b) Again traces exclude a minus sign in $HDH^{-1}=\pm D^{-1}$. Comparing entries in $HD=D^{-1}H$ forces both diagonal entries of $H$ to be zero. Its determinant then gives
\[
H=\begin{pmatrix}0&-r\\r^{-1}&0\end{pmatrix},\qquad r\ne0.
\]
Modulo sign, take $r>0$. These are $a_tJ$ with $r=e^{t/2}$ and $J=\bigl(\begin{smallmatrix}0&-1\\1&0\end{smallmatrix}\bigr)$.
(c) If $g=q a_\ell q^{-1}$, conjugate these answers by $q$. The commuting maps preserve $q(0)$ and $q(\infty)$ individually and translate along the same axis. The maps in (b) exchange the endpoints and reverse the translation along the axis. They still preserve orientation in the two-dimensional plane.
\end{solution}

\begin{solution}{ex:v9-ch09-direction-measure}
\idea An isometry changes the base point by its area Jacobian and adds a base-dependent angle to the direction.
\writeup (a) The transformed vector is $g'(z)y e^{i\alpha}$, whose Euclidean length is $\operatorname{Im}g(z)$. Its angle is locally $\alpha+\arg g'(z)$. In the full coordinate Jacobian the base point is independent of $\alpha$, while the angle derivative with respect to $\alpha$ is $1$. Its determinant is therefore the base-area Jacobian $|g'(z)|^2$. This preserves $dx\,dy/y^2$ and leaves $d\alpha$ as the fibre factor. Local angle branches suffice because their changes are multiples of $2\pi$. Under the smooth bijection $\Psi$, established by the Iwasawa coordinates, this is a left-invariant volume, hence Haar volume.
(b) In the basis $E,H,F$, $\operatorname{Ad}(a_t)$ has eigenvalues $e^t,1,e^{-t}$ and determinant $1$. Direct multiplication gives
\[
\operatorname{Ad}(n_x)E=E,\quad
\operatorname{Ad}(n_x)H=H-2xE,\quad
\operatorname{Ad}(n_x)F=F+xH-x^2E,
\]
so its triangular matrix also has determinant $1$. Conjugation by $k\in K$ preserves the positive inner product $\operatorname{tr}(X^TY)$ on real matrices, hence has determinant $\pm1$ on the invariant subspace $\mathfrak{sl}_2(\R)$. Continuity along the circle $K$ makes that determinant $1$. Iwasawa decomposition now gives determinant $1$ for every group element. The criterion of Chapter~8 proves unimodularity, so Haar volume is also invariant under right multiplication, including both flows.
(c) Almost every point of $F$ has no nontrivial stabilizer, and its direction fibre has length $2\pi$. The exceptional boundary and cone points have volume zero. Total volume is therefore $(\pi/3)2\pi=2\pi^2/3$. The specified event has volume $(1/h)\beta$, hence probability
$3\beta/(2\pi^2h)$. The formula includes $\beta=0$ and $2\pi$; it requires no claim about individual trajectories.
\end{solution}

\section*{Sources and further reading}
\addcontentsline{toc}{section}{Sources and further reading}

H.~Poincaré, \emph{Science et méthode} (Flammarion, 1908), Book~I, Chapter~3, ``L'invention mathématique''; English translation by G.~B.~Halsted in \emph{The Foundations of Science} (Science
Press, 1913). J.~Gray, \emph{Henri Poincaré: A Scientific Biography} (Princeton University Press, 2013), on the Fuchsian functions of 1880--81 and the Coutances episode.

C.~Series, ``The modular surface and continued fractions'', \emph{J.\ London Math.\ Soc.} (2) 31 (1985), 69--80; E.~Artin, ``Ein mechanisches System mit quasiergodischen Bahnen'', \emph{Abh.\ Math.\
Sem.\ Univ.\ Hamburg} 3 (1924), 170--175; E.~Hopf, ``Fuchsian groups and ergodic theory'', \emph{Trans.\ Amer.\ Math.\ Soc.} 39 (1936), 299--314; A.~Hurwitz, ``\"Uber algebraische Gebilde mit
eindeutigen Transformationen in sich'', \emph{Math.\ Ann.} 41 (1893), 403--442. Gauss's problem on continued fractions is in his letter to Laplace of 30 January 1812, \emph{Werke} X.1 (1917), 371--374.

The treatment follows S.~Katok, \emph{Fuchsian Groups} (University of Chicago Press, 1992), Chapters~1--4; M.~Einsiedler and T.~Ward, \emph{Ergodic Theory with a view towards Number Theory}
(Springer, 2011), Chapter~9, for the geodesic flow as a matrix product and the link with continued fractions; and C.~Series, ``Continued fractions and hyperbolic geometry'', lecture notes of
the London Mathematical Society summer school, Loughborough, 2015, which draws the Farey pictures beautifully. For Klein's quartic, S.~Levy (ed.), \emph{The Eightfold Way} (MSRI
Publications 35, Cambridge University Press, 1999).
